% !TEX root = ../main.tex
% ---------------------------------------------------------------------
%  P1 -- Virtual Memory / MMU
% ---------------------------------------------------------------------
\vlsec{cP1}{P1}{Virtual Memory \& MMU}
\label{sec:mmu}

\begin{ztab}{@{}lX@{}}
Cortex-M & MCU, $<$200\,MHz, meist kein Cache, Code aus Flash \\
Cortex-R & \T{MPU}: übersetzt \emph{nicht}, prüft nur Rechte (r/w/x) für
  typ.\ 8 Regionen \ra Taskisolation, „no execute“, Zugriffszeit bleibt
  vorhersagbar \\
Cortex-A & \T{MMU}: kann alles der MPU \emph{und} übersetzt Adressen.
  Voraussetzung für Linux. \\
\end{ztab}

\subsection{Drei Probleme, eine Lösung}
\begin{enumerate}
\item zu wenig RAM \item keine Isolation zwischen Prozessen
\item Fragmentierung des freien Speichers
\end{enumerate}
\SW{Idee} Jeder Prozess bekommt die \emph{Illusion} eines eigenen,
zusammenhängenden Adressraums. \T{Pages} (typ.\ 4\,kB) virtuell,
\T{Frames} physisch, \T{Page Table} ordnet zu.
\begin{itemize}
\item Aufeinanderfolgende Pages dürfen in verstreuten Frames liegen
      \ra Fragmentierung erledigt
\item Jeder Prozess hat eigene Tabelle mit disjunkten Frames
      \ra Isolation erledigt
\end{itemize}

\SW{Paging/Swapping} Reicht das RAM nicht, wandert eine selten benutzte
Page auf die Disk, ihr Eintrag wird invalid. Späterer Zugriff
\ra \T{Page Fault} \ra OS holt zurück, verdrängt eine andere, Befehl wird
wiederholt. Dauert $\sim$100\,000 Cycles
\ra \emph{page miss $=$ task switch}.

\subsection{Adressübersetzung}
\SW{Einstufig} 32\,Bit $=$ 20\,Bit Page-Nummer $+$ 12\,Bit Offset
($2^{12}=4$\,kB). Der Offset wird \emph{nicht} übersetzt (Pages und Frames
sind gleich gross). Page-Nummer indiziert die Tabelle, diese liefert
Frame-Nummer $+$ Rechte.

\begin{falle}[Warum einstufig nicht geht]
$2^{20}$ Einträge $\times$ 4\,B $=$ \T{4\,MB Page Table pro Prozess}.
Bei 50 Prozessen 200\,MB nur für Tabellen.
\end{falle}

\SW{Zweistufig} $10+10+12$. Die ersten 10\,Bit indizieren eine
First-Level-Tabelle (1024 Einträge, 4\,kB); jeder Eintrag zeigt auf eine
Second-Level-Tabelle mit wieder 1024 Einträgen. Ein Prozess, der nur
10\,MB benutzt, braucht wenige Second-Level-Tabellen \ra statt 4\,MB nur
ca.\ 20\,kB. Preis: ein zusätzlicher Speicherzugriff pro Übersetzung.

\subsection{TLB (Translation Lookaside Buffer)}
Ohne TLB kostete jeder Speicherzugriff zwei zusätzliche RAM-Zugriffe nur
für die Übersetzung -- dreimal langsamer als ganz ohne MMU. Der TLB ist
der Cache für Übersetzungen: klein, aber hoch assoziativ.

\SW{Warum klein reicht} Ein Eintrag deckt 4\,kB Adressraum ab; 32 Einträge
schon 128\,kB. Mit Locality sitzt ein Programm fast immer in wenigen Pages.

\begin{ztab}{@{}lXl@{}}
Cortex-A72 & Struktur & Einträge \\ \midrule
Micro-TLB D & voll assoziativ & 32 \\
Micro-TLB I & voll assoziativ & 48 \\
L2-TLB & 4-way set-assoziativ & 1024 \\
\end{ztab}

\begin{formel}[Die Analogie, die alles zusammenhält]
\begin{tabularx}{\linewidth}{@{}XX@{}}
\T{Cache} & \T{Virtual Memory} \\
Cache Line & Page / Frame \\
Tag-RAM & Page Table \\
Valid Bit & Valid Bit \\
Offset in der Line & Offset in der Page \\
Miss \ra nachladen & Page Fault \ra einlagern \\
Replacement in HW & Replacement im OS (SW) \\
\end{tabularx}\\[2pt]
Nur die Zeitskalen unterscheiden sich um Faktor $\sim$1000 -- deshalb
entscheidet dort Hardware, hier Software.
\end{formel}

\subsection{\R{vm\_stat} lesen}
\begin{ztab}{@{}lX@{}}
free/active/inactive & RAM-Zustand, Auslagerungskandidaten \\
wired down & nie auslagerbar (Kernel -- wer behandelte sonst den Fault?) \\
translation faults & Page Faults, meist harmlose \emph{soft faults} \\
pageins / pageouts & Disk-Transfers \\
swapins / swapouts & stetig steigend $=$ zu wenig RAM, System thrasht \\
copy-on-write & Fork teilt Frames bis zum ersten Schreibzugriff \\
\end{ztab}
Achtung: Apple Silicon benutzt 16-kB-Pages, nicht 4\,kB wie im Rest der
Vorlesung.
